\chapter{Breuken}
\label{chap:breuken}

Breuken komen we in het dagelijkse leven veel tegen: de helft van een 
pizza, drie kwartier wachttijd, of twee derde van de studenten waren op tijd voor de les. Voorbeelden van breuken zijn: 
\[
    \frac{1}{3}, \quad \frac{8}{15}, \quad \frac{2}{p}, \quad \frac{x}{y}
\]
Formeel beschreven is een breuk een getal met de volgende structuur: 
$\frac{T}{N}$. Hierbij noemen we $T$ de \emph{teller} van de breuk, 
terwijl we $N$ de \emph{noemer} van de breuk noemen.

In principe is een breuk een deling: je deelt de teller $T$ door de 
noemer $N$. Op deze wijze kun je de decimale waarde van een breuk 
berekenen. Bijvoorbeeld: 
\[
    \frac{8}{15}\approx 0{,}5333
\]

Goed om te benoemen is dat er \'e\'en belangrijke beperking is: \emph{de noemer mag nooit nul zijn}. Delen door nul is niet gedefinieerd, dus bij elke breuk (en elke formule met een deling!) moet je je afvragen of de noemer wel ongelijk aan nul is.

Waarom is het rekenen met breuken zo belangrijk? Veel studenten proberen 
breuken te omzeilen door alles met de rekenmachine in decimalen om te 
zetten. Maar vaak is een breuk beter, een breuk als $\frac{2}{3}$ is namelijk \emph{exact}, terwijl de decimale notatie $0.6667$ afgerond (en dus onnauwkeuriger) is.

In dit hoofdstuk behandelen we een aantal eigenschappen van breuken: 
het vereenvoudigen en gelijknamig maken van breuken, en het optellen, 
aftrekken, vermenigvuldigen en delen van breuken. Ook dit gaat weer aan de hand van enkele rekenregels. 

\newpage
\regel{Regel 1}{\frac{A}{B}=\frac{p\cdot A}{p\cdot B} \quad \text{waarbij } p\neq 0}

In woorden: de waarde van een breuk verandert niet, als we de teller en de noemer met hetzelfde getal vermenigvuldigen.

Voorbeelden:
\begin{align*}
\frac{8}{15}    &= \frac{2 \cdot 8}{2 \cdot 15} = \frac{16}{30} & &\text{De breuken } \frac{8}{15} \text{ en } \frac{16}{30} \text{ hebben dezelfde waarde.} \\[1em]
\frac{5}{y}     &= \frac{3 \cdot 5}{3 \cdot y} = \frac{15}{3y} & &\text{De breuken } \frac{5}{y} \text{ en } \frac{15}{3y} \text{ hebben dezelfde waarde.}
\end{align*}

\regel{Regel 2}{\frac{A}{B} = \frac{\frac{A}{q}}{\frac{B}{q}} \quad \text{waarbij } q \neq 0}

In woorden: de waarde van een breuk verandert niet, als we de teller en de noemer door hetzelfde getal delen.

Voorbeelden:
\begin{align*}
\frac{10}{24} &= \frac{\frac{10}{2}}{\frac{24}{2}} = \frac{5}{12} & &\text{De breuken } \frac{10}{24} \text{ en } \frac{5}{12} \text{ hebben dezelfde waarde.} \\[1em]
\frac{24}{16x} &= \frac{\frac{24}{8}}{\frac{16x}{8}} = \frac{3}{2x} & &\text{De breuken } \frac{24}{16x} \text{ en } \frac{3}{2x} \text{ hebben dezelfde waarde.}
\end{align*}

\regel{Regel 3}{\frac{-A}{B}=\frac{A}{-B}=-\frac{A}{B} }

Voorbeelden:
\begin{align*}
\frac{-8}{15} &= -\frac{8}{15} \\[1em]
\frac{5}{-y}  &= -\frac{5}{y}  \\[1em]
-\frac{3}{a}  &= \frac{-3}{a}
\end{align*}

\section{Het vereenvoudigen van een breuken met een breuk in de breuk}
Er zijn regels waarmee je een breuk met een breuk in de teller en/of noemer kunt vereenvoudigen.

\regel{Regel 4}{\frac{\frac{A}{B}}{C}=\frac{A}{B\cdot C}}

Voorbeelden:
\begin{align*}
\frac{\frac{2}{5}}{3}    &= \frac{2}{5 \cdot 3} = \frac{2}{15} \\[1em]
\frac{\frac{3}{2}}{7}    &= \\[1em]
\frac{\frac{3}{x}}{y}    &= \frac{3}{x \cdot y} = \frac{3}{xy} \\[1em]
\frac{\frac{5a}{3b}}{4c} &= \\[1em]
\frac{\frac{2}{x}}{3x}   &= \frac{2}{x \cdot 3x} = \frac{2}{3x^2} \\[1em]
\frac{\frac{y}{5}}{y}    &=
\end{align*}

\regel{Regel 5}{\frac{A}{\frac{B}{C}}=\frac{A\cdot C}{B}}

Voorbeelden:
\begin{align*}
\frac{2}{\frac{5}{3}}    &= \frac{2 \cdot 3}{5} = \frac{6}{5} \\[1em]
\frac{3}{\frac{2}{7}}    &= \\[1em]
\frac{3}{\frac{x}{2y}}   &= \frac{3 \cdot 2y}{x} = \frac{6y}{x} \\[1em]
\frac{5a}{\frac{3b}{4c}} &= \\[1em]
\frac{2}{\frac{x}{3x}}   &= \frac{2 \cdot 3x}{x} = \frac{6x}{x} = 6 \\[1em]
\frac{y}{\frac{5}{15y}}  &=
\end{align*}

\newpage
\regel{Regel 6}{\frac{\frac{A}{B}}{\frac{C}{D}}=\frac{A\cdot D}{B\cdot C}}

Voorbeelden:
\begin{align*}
\frac{\frac{4}{5}}{\frac{3}{7}}      &= \frac{4 \cdot 7}{5 \cdot 3} = \frac{28}{15} \\[1em]
\frac{\frac{7}{3}}{\frac{2}{5}}      &= \\[1em]
\frac{\frac{2a}{3m}}{\frac{6b}{5k}}  &= \frac{2a \cdot 5k}{3m \cdot 6b} = \frac{10ak}{18mb} = \frac{5ak}{9mb} \\[1em]
\frac{\frac{8c}{5r}}{\frac{4d}{10s}} &= \\[1em]
\frac{\frac{x}{y}}{\frac{y}{x}}      &= \frac{x^2}{y^2} \\[1em]
\frac{\frac{d}{s}}{\frac{d^2}{s^2}}  &=
\end{align*}

\section{Het vermenigvuldigen van breuken}
Wanneer breuken met elkaar worden vermenigvuldigd, dan kunnen we direct de tellers met elkaar vermenigvuldigen en ook de noemers met elkaar vermenigvuldigen (zie onderstaande regel 7). We hoeven dus niet eerst de breuken gelijknamig te maken zoals bij het optellen van breuken (zie paragraaf 5.3).

\regel{Regel 7}{\frac{A}{B}\cdot \frac{C}{D}=\frac{A\cdot C}{B\cdot D}}

Voorbeelden:
\begin{align*}
\frac{3}{5} \cdot \frac{2}{7}     &= \frac{3 \cdot 2}{5 \cdot 7} = \frac{6}{35} \\[1em]
\frac{4}{9} \cdot \frac{3}{5}     &= \\[1em]
\frac{3}{x} \cdot \frac{2}{y}     &= \frac{6}{xy} \\[1em]
\frac{a}{2} \cdot \frac{3}{b}     &= \\[1em]
x \cdot \frac{y}{3}               &= \frac{x}{1} \cdot \frac{y}{3} = \frac{xy}{3} \\[1em]
c \cdot \frac{a}{b}               &= \\[1em]
\frac{2+x}{3-y} \cdot \frac{5}{8} &= \frac{(2+x) \cdot 5}{(3-y) \cdot 8} = \frac{10+5x}{24-8y} \\[1em]
\frac{a+2}{b-1} \cdot \frac{b}{a} &= \\[1em]
\frac{x-1}{x-2} \cdot \frac{x-2}{x-3} &= \frac{(x-1)(x-2)}{(x-2)(x-3)} = \frac{x-1}{x-3} \\[1em]
\frac{y+4}{y-1} \cdot \frac{y+3}{y+4} &= \\[1em]
\frac{2}{a} \cdot \frac{b^2}{4} \cdot \frac{8}{c}    &= \frac{16b^2}{4ac} = \frac{4b^2}{ac} \\[1em]
\frac{x}{2y} \cdot \frac{10}{p} \cdot \frac{y^2}{20} &=
\end{align*}

\section{Het optellen van breuken}
Wanneer breuken bij elkaar opgeteld worden, moeten ze eerst \emph{gelijknamig} gemaakt worden, d.w.z. we moeten ervoor zorgen dat ze gelijke noemers krijgen. Uiteraard hoeft dat alleen maar, als ze nog niet gelijknamig zijn. Als we de gelijknamigheid bereikt hebben, kunnen we de breuken bij elkaar optellen door de tellers van deze gelijknamige breuken bij elkaar op te tellen (zie regel 8).

\regel{Regel 8}{\frac{A}{N}+\frac{B}{N}=\frac{A+B}{N}}

Voorbeelden:
\begin{align*}
\frac{3}{7} + \frac{5}{7}       &= \frac{8}{7} \\[1em]
\frac{5}{12} + \frac{7}{12}     &= \\[1em]
\frac{4}{x} + \frac{5}{x}       &= \frac{4+5}{x} = \frac{9}{x} \\[1em]
\frac{6}{p} + \frac{2}{p}       &= \\[1em]
\frac{4}{5+a} + \frac{2}{5+a}   &= \frac{4+2}{5+a} = \frac{6}{5+a} \\[1em]
\frac{5}{x-2} + \frac{6}{x-2}   &= \\[1em]
\frac{3}{7} + \frac{1}{2}       &= \frac{6}{14} + \frac{7}{14} = \frac{13}{14} \\[1em]
\frac{1}{3} + \frac{2}{5}       &= \\[1em]
\frac{8}{x} - \frac{5}{y}       &= \frac{8y}{xy} - \frac{5x}{xy} = \frac{8y-5x}{xy} \\[1em]
\frac{4}{a} + \frac{2}{b}       &= \\[1em]
\frac{3}{p+2} + \frac{5}{p}     &= \frac{3p}{p(p+2)} + \frac{5(p+2)}{p(p+2)} = \frac{3p}{p(p+2)} + \frac{5p+10}{p(p+2)} = \frac{3p+5p+10}{p(p+2)} = \frac{8p+10}{p(p+2)} \\[1em]
\frac{5}{x} - \frac{3}{x+4}     &= \\[1em]
\frac{3a}{a-1} + \frac{4}{a-3}  &= \\[1em]
\frac{1}{2} + \frac{5}{6}       &= \frac{3}{6} + \frac{5}{6} = \frac{8}{6} = \frac{4}{3} \\[1em]
\frac{3}{10} + \frac{2}{15}     &= \\[1em]
\frac{x}{y} - \frac{x}{3y}      &= \frac{3x}{3y} - \frac{x}{3y} = \frac{3x-x}{3y} = \frac{2x}{3y} \\[1em]
\frac{2}{x} - \frac{1}{x^2}     &= \\[1em]
\frac{1}{2} + \frac{2}{3} - \frac{3}{5} &= \frac{15}{30} + \frac{20}{30} - \frac{18}{30} = \frac{17}{30} \\[1em]
\frac{5}{a} - \frac{4}{b} + \frac{3}{c} &=
\end{align*}

\section{Een breuk vereenvoudigen}
Soms zijn er breuken waarbij in de teller of noemer een breuk voorkomt als onderdeel van die teller of noemer. Het doel bij het vereenvoudigen hiervan is om de breuk, die in de teller of de noemer staat, weg te werken.

Voorbeelden:
\begin{align*}
\frac{3+\frac{2}{5}}{4} &\quad \text{Dit vraagstuk gaan we op 2 manieren aanpakken:} \\[0.5em]
\frac{3+\frac{2}{5}}{4} &= \frac{\frac{15}{5}+\frac{2}{5}}{4} = \frac{\frac{17}{5}}{4} = \frac{17}{5 \cdot 4} = \frac{17}{20} & &\text{Manier 1} \\[1em]
\frac{3+\frac{2}{5}}{4} &= \frac{3+\frac{2}{5}}{4} \cdot \frac{5}{5} = \frac{(3+\frac{2}{5}) \cdot 5}{4 \cdot 5} = \frac{15+2}{20} = \frac{17}{20} & &\text{Manier 2} \\[1em]
\frac{5-\frac{2}{3}}{5} &= \\[1em]
\frac{1+\frac{2}{x}}{y} &= \frac{1+\frac{2}{x}}{y} \cdot \frac{x}{x} = \frac{(1+\frac{2}{x}) \cdot x}{y \cdot x} = \frac{x+2}{xy} & &\text{Hier is de 2e manier toegepast} \\[1em]
\frac{3-\frac{2}{a}}{b} &= \\[1em]
\frac{4}{1-\frac{1}{a}} &= \frac{4}{\frac{a}{a}-\frac{1}{a}} = \frac{4}{\frac{a-1}{a}} = \frac{4a}{a-1} & &\text{Hier is de 1e manier toegepast} \\[1em]
\frac{4c}{\frac{1}{d}-2} &= \\[1em]
\frac{1-\frac{1}{k}}{3+\frac{2}{k}} &= \frac{1-\frac{1}{k}}{3+\frac{2}{k}} \cdot \frac{k}{k} = \frac{(1-\frac{1}{k}) \cdot k}{(3+\frac{2}{k}) \cdot k} = \frac{k-1}{3k+2} & &\text{Hier is de 2e manier toegepast} \\[1em]
\frac{5+\frac{2}{m}}{1-\frac{1}{m}} &= \\[1em]
\frac{1+\frac{3}{a}}{2+\frac{5}{b}} &= \frac{\frac{a}{a}+\frac{3}{a}}{\frac{2b}{b}+\frac{5}{b}} = \frac{\frac{a+3}{a}}{\frac{2b+5}{b}} = \frac{(a+3) \cdot b}{a \cdot (2b+5)} = \frac{ab+3b}{2ab+5a} & &\text{Hier is de 1e manier toegepast} \\[1em]
\frac{\frac{1}{x}-4}{3-\frac{5}{y}} &= \\[1em]
\frac{\frac{2}{a}-\frac{4}{b}}{\frac{3}{b}+\frac{1}{a}} &= \frac{\frac{2b}{ab}-\frac{4a}{ab}}{\frac{3a}{ab}+\frac{b}{ab}} = \frac{\frac{2b-4a}{ab}}{\frac{3a+b}{ab}} = \frac{(2b-4a) \cdot ab}{ab \cdot (3a+b)} = \frac{2b-4a}{3a+b} & &\text{Hier is de 2e manier toegepast} \\[1em]
\frac{\frac{4}{c}+\frac{5}{d}}{\frac{3}{d}-\frac{1}{c}} &=
\end{align*}

\section{Huiswerkopdrachten}
In deze opgaven pas je de rekenregels met breuken toe die je in dit hoofdstuk hebt geleerd. Werk elke opgave netjes uit en vereenvoudig je eindantwoord zo ver mogelijk.

\subsection*{Opgave 1: Vereenvoudig}
\renewcommand{\arraystretch}{2}
\begin{tabular}{m{15em}@{\hspace{4em}}l}
\text{a.}\hspace{0.5em} $\frac{2y}{6x}$        & \text{j.}\hspace{0.5em} $\frac{a}{\frac{b}{ac}}$        \\
\text{b.}\hspace{0.5em} $\frac{5}{2}$         & \text{k.}\hspace{0.5em} $\frac{\frac{abc}{d}}{\frac{ac}{bd}}$  \\
\text{c.}\hspace{0.5em} $\frac{4y}{3x}$      & \text{l.}\hspace{0.5em} $\frac{\frac{2x}{y}}{6xy}$      \\
\text{d.}\hspace{0.5em} $\frac{x}{\frac{3y}{4}}$  & \text{m.}\hspace{0.5em} $\frac{2x}{6xy}$   \\
\text{e.}\hspace{0.5em} $\frac{3}{\frac{2}{x}}$  & \text{n.}\hspace{0.5em} $\frac{\frac{bc}{a}}{abc}$  \\
\text{f.}\hspace{0.5em} $\frac{3xy}{3y}$      & \text{o.}\hspace{0.5em} $\frac{abc}{\frac{bc}{a}}$  \\
\text{g.}\hspace{0.5em} $\frac{8a}{2ab}$      & \text{p.}\hspace{0.5em} $\frac{6km}{\frac{km}{6}}$  \\
\text{h.}\hspace{0.5em} $\frac{a}{2ab}$       & \text{q.}\hspace{0.5em} $\frac{\frac{ab}{a}}{\frac{a}{bc}}$  \\
\text{i.}\hspace{0.5em} $\frac{2x}{\frac{x}{2}}$  & \text{r.}\hspace{0.5em} $\frac{\frac{8xy}{z}}{\frac{8y}{xz}}$  \\
\end{tabular}

\subsection*{Opgave 2: Herleid tot 1 breuk en vereenvoudig daarna indien mogelijk}
\renewcommand{\arraystretch}{2}
\begin{tabular}{m{15em}@{\hspace{4em}}l}
\text{a.}\hspace{0.5em} $\frac{4}{x} \cdot \frac{p}{y}$           & \text{k.}\hspace{0.5em} $\frac{2-k}{k+4} \cdot \frac{k-3}{2-k}$  \\
\text{b.}\hspace{0.5em} $\frac{a}{x} \cdot \frac{5}{y}$           & \text{l.}\hspace{0.5em} $\frac{a+2}{b-1} \cdot \frac{b}{a}$     \\
\text{c.}\hspace{0.5em} $k \cdot \frac{2k}{m}$                   & \text{m.}\hspace{0.5em} $\frac{8x^3}{y^5} \cdot \frac{y^2}{4x}$ \\
\text{d.}\hspace{0.5em} $6 \cdot \frac{2x-1}{3}$                 & \text{n.}\hspace{0.5em} $2p \cdot \frac{6}{p-6}$               \\
\text{e.}\hspace{0.5em} $\frac{a}{b} \cdot \frac{b}{2a}$         & \text{o.}\hspace{0.5em} $\frac{2k}{m} \cdot \frac{k+4}{5-m}$   \\
\text{f.}\hspace{0.5em} $\frac{a^2}{b^3} \cdot \frac{b^2}{a}$    & \text{p.}\hspace{0.5em} $\frac{ac}{b^3} \cdot \frac{b^2}{c} \cdot \frac{bc^2}{a}$ \\
\text{g.}\hspace{0.5em} $3 \cdot \frac{2-x}{2y}$                 & \text{q.}\hspace{0.5em} $\frac{a}{b^2} \cdot \frac{c}{b} \cdot \frac{ab}{c}$ \\
\text{h.}\hspace{0.5em} $a \cdot \frac{3}{x+2}$                  & \text{r.}\hspace{0.5em} $2x \cdot \frac{5}{3y^2} \cdot \frac{6y}{x}$ \\
\text{i.}\hspace{0.5em} $\frac{2-x}{3} \cdot \frac{3}{2}$        & \\
\text{j.}\hspace{0.5em} $\frac{a+5}{a-1} \cdot \frac{a-3}{a+5}$  & \\
\end{tabular}

\subsection*{Opgave 3: Herleid tot 1 breuk en vereenvoudig}
\renewcommand{\arraystretch}{2}
\begin{tabular}{m{15em}@{\hspace{4em}}l}
\text{a.}\hspace{0.5em} $\frac{2}{x} - \frac{3}{x} - \frac{5}{x}$   & \text{j.}\hspace{0.5em} $\frac{7}{3k} - \frac{3}{2k}$         \\
\text{b.}\hspace{0.5em} $-\frac{3}{y} + \frac{5}{y} - \frac{1}{y}$ & \text{k.}\hspace{0.5em} $\frac{a}{a+b} + \frac{a}{a-b}$     \\
\text{c.}\hspace{0.5em} $\frac{4a}{x+5} + \frac{3a}{x+5}$          & \text{l.}\hspace{0.5em} $\frac{2}{y} - \frac{4}{y-3}$       \\
\text{d.}\hspace{0.5em} $\frac{5x}{x-3} - \frac{2x}{x-3}$          & \text{m.}\hspace{0.5em} $\frac{m}{m-2} - \frac{m}{m+2}$     \\
\text{e.}\hspace{0.5em} $x - \frac{x}{x+1}$                        & \text{n.}\hspace{0.5em} $\frac{1}{x} + \frac{1}{x^2} - \frac{1}{x^3}$ \\
\text{f.}\hspace{0.5em} $3 - \frac{a}{b}$                          & \text{o.}\hspace{0.5em} $\frac{a}{b} - \frac{a}{c} + \frac{a}{d}$ \\
\text{g.}\hspace{0.5em} $b + \frac{a}{b}$                          & \text{p.}\hspace{0.5em} $\frac{a+c}{c^2} - \frac{1}{c}$     \\
\text{h.}\hspace{0.5em} $\frac{5}{x-2} - \frac{3}{x-1}$            & \text{q.}\hspace{0.5em} $\frac{5x}{x+1} + \frac{2}{x+2}$    \\
\text{i.}\hspace{0.5em} $\frac{7a}{5b} - \frac{3a}{b}$            & \\
\end{tabular}

\subsection*{Opgave 4: Vereenvoudig}
\renewcommand{\arraystretch}{2}
\begin{tabular}{m{15em}@{\hspace{4em}}l}
\text{a.}\hspace{0.5em} $\frac{2+\frac{1}{x}}{x}$              & \text{j.}\hspace{0.5em} $\frac{c-\frac{4}{d}}{c-\frac{1}{d}}$        \\
\text{b.}\hspace{0.5em} $\frac{5-\frac{2}{3}}{2+\frac{1}{6}}$  & \text{k.}\hspace{0.5em} $\frac{\frac{a}{b}+2}{\frac{a}{b}-2}$          \\
\text{c.}\hspace{0.5em} $\frac{2}{1+\frac{a}{b}}$              & \text{l.}\hspace{0.5em} $\frac{\frac{y}{x}}{1-\frac{y}{x}}$           \\
\text{d.}\hspace{0.5em} $\frac{3+\frac{2}{p}}{p}$              & \text{m.}\hspace{0.5em} $\frac{1+\frac{2}{a}}{3+\frac{1}{b}}$         \\
\text{e.}\hspace{0.5em} $\frac{2+\frac{x}{y}}{1-x}$            & \text{n.}\hspace{0.5em} $\frac{\frac{m}{k}-4}{\frac{k}{m}+4}$         \\
\text{f.}\hspace{0.5em} $\frac{a+1}{1-\frac{b}{a}}$            & \text{o.}\hspace{0.5em} $\frac{2-\frac{4}{c}}{\frac{1}{b}-1}$         \\
\text{g.}\hspace{0.5em} $\frac{2+\frac{1}{x}}{\frac{5}{x}-1}$ & \text{p.}\hspace{0.5em} $\frac{1+\frac{3}{x}}{2+\frac{5}{y}}$         \\
\text{h.}\hspace{0.5em} $\frac{3k-\frac{1}{k}}{\frac{1}{k}+2}$ & \text{q.}\hspace{0.5em} $\frac{\frac{2}{c}-\frac{3}{d}}{4+\frac{1}{c}}$ \\
\text{i.}\hspace{0.5em} $\frac{x+\frac{1}{y}}{x-\frac{1}{y}}$  & \text{r.}\hspace{0.5em} $\frac{\frac{2}{k}-\frac{5}{m}}{1+\frac{3}{k}}$ \\
\end{tabular}
